\head{第二十四版整合核查}
\label{audit:V24-integration}

\audit{V24-I01}{实际底本、范围与下一接续点}{本轮以已发布第二十三版的实际工程为基线，逐段对读法文与英文载体。}
本轮范围是第二篇第16节“函数的奇异零点与超奇异零点”：法文扫描物理95页起至111页最后原注；英文编订稿物理62页节题起至73页第17节标题之前。法文物理112页、英文物理73页中部开始第17节，属于下一回校范围。短篇\texttt{EGAV1-2.ps}的同文题名、编者归属及内容与本节对应处理；同文在正文只排印一次。具体逐段来源、公式和交换图的读法分别见本版两份专题报告。

\audit{V24-I02}{实际工程恢复与历史字节保存}{先核对已发布文件，再建立本版副本；历史记录不代替当轮检查。}
本轮开始时，前轮暂存工程已不在工作目录中；现有工程经公开仓库的普通克隆及显式检出恢复。基线提交为\texttt{aaa56a39\allowbreak 4186471c\allowbreak 071e9d9c\allowbreak c686401b\allowbreak ecadf98e}。恢复后实际读取并核对仓库清单：7869项载荷文件、总计1867769118字节，严格校验没有报告差异，工作树干净。此处记载的是本轮起点的实际恢复检查，不是本版最终清单的文件数。

第二十三版当前工程共1285项文件、231701353字节，已逐项比较后复制为本版工作副本。需要更新的入口、版次说明及当前检查记录，其第二十三版字节另存于\path{recollation/legacy_v23/}；其余历史内容沿用原路径。保存表记录每项原路径、保存路径、字节数及校验值。四个正文、续作、报告和书末资料目录中的207项既有TeX文件均纳入继承清单。上述保存和编排检查只证明历史文件未被漏掉或改写，不表示这些文件中的所有数学判断已在本轮重新审定。

\audit{V24-I03}{第16节替换及Dot保留}{新模块在原位置接续，旧全文完整保存，避免两份同文同时进入合订本。}
既有\path{text/v1_2.tex}包含第16节译文及后部历史译校附记，共61675字节。本版将它作为停用原始资料保存，不在正文重复输入；两份新模块\path{continuations/v24_section16_early.tex}与\path{continuations/v24_section16_late.tex}共同替代其唯一活动输入。旧附记中与本节有关的内容按可核证的适用范围吸收进脚注和专题报告，不以历史批次标题重复排印正文。

Dot第69--108批的40个活动模块继续各输入一次。第一篇第15节的第二十三版回校模块及此前各卷、前版报告与书末资料，按原有内容次序保留。源码检查分别比较活动输入顺序、既有字节、当前单文件展开及停用文件状态；本节完整替换不改变此前各卷正文的字节。

\audit{V24-I04}{订正、原错脚注与核查边界}{把有证据的订正放入正文，把原写法及其问题留在脚注和报告。}
本轮沿用使用者确定的译校原则：对于已确认的原文错误，正文写出有依据的正确版本，并以脚注注明原写法、错误所在及订正理由；原稿缺文、未决的一般化和不可辨内容保留标识。原作者论证、原作者删划、英文编者说明、旧译者补证与本轮译者补证须分别标明。对某一附加假设下的证明，不解释为无条件证明了原断言。本版专题报告记录具体原文位置、问题归属、反例或补证及尚存边界。

\audit{V24-I05}{出稿核验的含义}{实际文件检查、编译、逐页机器比对和指定页面目视核验各自记录，不能互相替代。}
正文与报告使用真实\texttt{ctexbook}类的默认字体和正文版式，公式及交换图以TeX排版。工程的出稿记录分别绑定活动源码、模块版与单文件版、四份最终编译日志及对应PDF；逐页机器比对覆盖文本、页面尺寸和72\,dpi像素。实际目视检查另列所看的物理页和图像，不能被解释为合订本全部历史页面已逐页人工复审。

最终生成的页数、日志结果、逐页比对结果、实际目视页集合和文件校验值，以工程\path{ARTIFACTS.json}、\path{HANDOVER_STATE.json}及\path{checks/}中的当版记录为准。报告不以预计页数或未运行的检查填充结果。用户所提供全套材料的翻译与数学核查尚未完成；正式第IV章仍需从原卷内容提要、章首说明及（1.1.1）接续，本版预备手稿第16节回校不替代该项工作。
